\clearpage
\section{Gravity measurements}
\begin{colbox}
    \begin{enumerate}
        \item How many cycles will be required to measure \(g\) within 1\%?
        \item How many measurements wil be needed to verify the theory within \(1.65\) sigma?
        \item Calculate the error for \(g\) and write down an estimator for its error
        \item Estimate the average of gravity using the measurements, then estimate the variance using the measurements. Compare to the theoretical.
        \item Calculaute \(\hat{\sigma}_\sigma\) for parts 3 and 4, what causes the difference between these values?
    \end{enumerate}
\end{colbox}
\subsection{g within 1 percent}
First we reverse the relation
\begin{equation}
    T = 2\pi\sqrt\frac{l}{g} \implies g = \frac{4\pi^2l}{T^2}
\end{equation}
therefore the error in \(g\) is
\begin{equation}
    \Delta g = 4\pi^2l\sqrt{\frac{4}{T^6}\Delta T^2} = \frac{8\pi^2l}{T^3}\Delta T
\end{equation}
The number of cycles is such that
\begin{equation}
    T = \frac{t}{N} \implies \Delta T = \frac{\Delta t}{N}
\end{equation}
returning, we want
\begin{equation}
    \frac{\Delta g}{g} = \frac{\frac{8\pi^2l}{T^3}\Delta T}{\frac{4\pi^2l}{T^2}} = \frac{2}{T}\Delta T \leq 0.01
\end{equation}
inserting the cycles
\begin{equation}
    \frac{2\Delta t}{t} \leq 0.01
\end{equation}
inserting \(\Delta t = 0.5\)
\begin{equation}
    \frac{\qty{1}{\second}}{t}\leq 0.01 \implies t \geq \qty{100}{\second}
\end{equation}
the cycle is of length
\begin{equation}
    T = 2\pi\sqrt{\frac{l}{g}} \approx \qty{2}{\second}
\end{equation}
therefore
\begin{equation}
    N = \frac{t}{T} \geq \frac{\qty{100}{\second}}{\qty{2}{\second}} = 50
\end{equation}
\subsection{1.65 sigma}
The gravitational acceleration is
\begin{equation}
    g = \frac{GM}{r^2}
\end{equation}
so 
\begin{align}
    g_{sea} = \frac{GM}{r_{sea}^2} && g_{everest} = \frac{GM}{r_{everest}^2}
\end{align}
we can express the everest acceleration using the sea acceleration
\begin{equation}
    g_{everest} = g_{sea}\ins{\frac{r_{sea}}{r_{everest}}}^2
\end{equation}
we can insert the known numbers
\begin{equation}
    g_{everest} = \qty{9.8}{\meter\per\second^2}\cdot\ins{\frac{6400}{6409}}^2 = 9.7725
\end{equation}
we want 
\begin{equation}
    \frac{g_{sea}-g_{everest}}{\Delta g} = \frac{0.0275}{\Delta g} > 1.65 \implies \Delta g < \qty{0.01667}{\meter\per\second^2}
\end{equation}
but
\begin{equation}
    \Delta g = \frac{8\pi^2l}{T^3}\Delta T =  \frac{8\pi^2l}{T^3}\frac{\Delta t}{N}
\end{equation}
rearranging
\begin{equation}
    N = \frac{8\pi^2l}{T^3}\frac{\Delta t}{\Delta g}
\end{equation}
inserting all the numbers we find
\begin{equation}
    N = \frac{8\pi^2\cdot\qty{1}{m}}{\ins{\qty{2}{\second}}^3}\frac{\qty{0.5}{\second}}{\qty{0.01667}{\meter\per\second^2}} > 296
\end{equation}
\subsection{Student Experiment}
This time both the time and length have errors, therefore
\begin{equation}
    \Delta g  =\sqrt{\ins{\pdv{g}{l}\Delta l}^2+\ins{\pdv{g}{T}\Delta T}^2}
\end{equation}
where
\begin{equation}
    g = \frac{4\pi^2l}{T^2}
\end{equation}
giving us 
\begin{equation}
    \Delta g = \sqrt{\ins{\frac{4\pi^2}{T^2}\cdot\qty{0.07}{\meter}}^2+\ins{\frac{8\pi^2l}{T^3}\cdot\qty{0.05}{\second}}^2} \approx \qty{0.85}{\meter\per\second^2}
\end{equation}
but since we have 50 measurements the total standard deviation is
\begin{equation}
    \sigma = \frac{\Delta g}{\sqrt{50}} = \qty{0.12}{\meter\per\second^2}
\end{equation}
we can use the formula for an estimator for a variance with a known mean:
\begin{equation}
    \hat{\Delta g} = \sqrt{\frac{1}{N}\sum_i\ins{g_i-\frac{4\pi^2l_i}{T_i^2}}^2}
\end{equation}
\subsection{With Measurements}
A good estimate for the mean is simply the average of the measurements:
\begin{equation}
    \hat{g} = \overline{g} = \qty{9.792}{\meter\per\second^2}
\end{equation}
and the variance can be calculated using the formula for an unknown mean
\begin{equation}
    \hat{V} = \frac{1}{N-1}\sum_i\ins{g_i-\overline{g}}^2 = \qty{0.094}{(\meter\per\second^2)^2}
\end{equation}
therefore 
\begin{equation}
    \sigma = \qty{0.307}{\meter\per\second^2}
\end{equation}
\subsection{sigma sigma}
\begin{align}
    \hat{\sigma}_{\sigma,c} &= \frac{\sigma}{\sqrt{2N}} = \qty{0.012}{\meter\per\second^2} \\
    \hat{\sigma}_{\sigma,d} &= \frac{\sigma}{\sqrt{2N}} = \qty{0.0307}{\meter\per\second^2}
\end{align}
The sigma for section 3 is smaller than that of section 4. This is because in section 3 we took the value of \(g\) as an assumption for a value we know, while in section 4 we did not. This prior knowledge changes our uncertainty.